\subsection{Separable algebraic elements}

DEFINITION 3. --- Let E be an extension of K. An element x of E which is algebraic over K is said to be separable over K if the algebraic extension \(\mathrm{K}(x)\) of K is separable.

PROPOSITION 5. --- Let E be an extension of K, x an element of E algebraic over K and f the minimal polynomial of x over K. Then the following conditions are equivalent :

\begin{enumerate}
\def\labelenumi{\alph{enumi})}
\item
  \(x\) is separable over \(\mathbf{K}\) ;
\item
  the polynomial \(f\) is separable;
\item
  \(x\) is a simple root of \(f\) .
\end{enumerate}

The equivalence of \(a)\) and \(b)\) follows from Prop. 3, that of \(b)\) and \(c)\) from Prop. 3 and 4 (cf.~V, p.~37 and 38).

COROLLARY 1. --- If an element x of E is a simple root of a polynomial g of K{[}X{]}, it is separable over K.

For the minimal polynomial \(f\) of \(x\) over \(K\) divides \(g\) in \(K[X]\) (V, p.~16, Th. 1), so \(x\) is a simple root of \(f\) .

COROLLARY 2. --- If an element x of E is algebraic and separable over K, it is algebraic and separable over every extension K' of K contained in E.

Let \(f\) be the minimal polynomial of \(x\) over \(K\) . Then \(x\) is a simple root of \(f\) by Prop. 5, and since \(f\) belongs to \(K'[X]\) , the element \(x\) of \(E\) is separable over \(K'\) by Cor. 1.

COROLLARY 3. --- Suppose that K has characteristic \(p \neq 0\) . For an element \(x\) of E to belong to K it is necessary and sufficient for it to be both separable algebraic and \(p\) -radical over K.

The stated condition is clearly necessary. Conversely suppose that \(x\) is separable algebraic over \(\mathbf{K}\) and \(p\) -radical of height \(e\) over \(\mathbf{K}\) . Since \(x\) is separable over \(\mathbf{K}\) , the minimal polynomial \(f\) of \(x\) over \(\mathbf{K}\) does not belong to \(\mathbf{K}[X^p]\) (Prop. 4 and 5); since \(x\) is \(p\) -radical of height \(e\) over \(\mathbf{K}\) , we have \(f(\mathbf{X}) = \mathbf{X}^{p^e} - x^{p^e}(\mathbf{V},\text{p. 24, Prop. 1})\) ; we conclude that \(e = 0\) , so \(x \in \mathbf{K}\) .

PROPOSITION 6. --- Let E be an extension of K.

\begin{enumerate}
\def\labelenumi{\alph{enumi})}
\item
  If \(\mathbf{E}\) is algebraic and separable over \(\mathbf{K}\) , every element of \(\mathbf{E}\) is algebraic and separable over \(\mathbf{K}\) .
\item
  Conversely, let \(\mathbf{A}\) be a set of elements of \(\mathbf{E}\) , algebraic and separable over \(\mathbf{K}\) and such that \(\mathbf{E} = \mathbf{K}(\mathbf{A})\) ; then \(\mathbf{E}\) is algebraic and separable over \(\mathbf{K}\) .
\end{enumerate}

If \(\mathbf{E}\) is algebraic and separable over \(\mathbf{K}\) , the same is true of the extension \(\mathbf{K}(x)\) of \(\mathbf{K}\) for every \(x \in \mathbf{E}\) , whence \(a\) ).

Under the hypothesis \(b\) ), the extension \(\mathbf{E}\) is algebraic over \(\mathbf{K}\) (V, p.~18, Cor. 1).

Let \(\mathbf{F}\) be a subextension of \(\mathbf{E}\) of finite degree over \(\mathbf{K}\) . By \(\mathbf{V}\) , p.~11, Cor., there exist elements \(x_{1},\ldots ,x_{m}\) of \(\mathbf{A}\) such that \(\mathbf{F}\subset \mathbf{K}(x_1,\dots,x_m)\) and we have

\[
\mathrm{K} \left(x _ {1}, \dots , x _ {m}\right) = \mathrm{K} \left[ x _ {1}, \dots , x _ {m} \right] \quad (\mathrm{V}, \text { p. } 1 8, \text { Cor. } 1).
\]

By the hypothesis on A, the algebras \(K[x_{1}], \ldots, K[x_{m}]\) are etale over K; hence the same is true of \(K[x_{1}] \otimes \cdots \otimes K[x_{m}]\) (V, p.~32, Cor. 1). Now F is isomorphic to a subalgebra of a quotient algebra of \(K[x_{1}] \otimes \cdots \otimes K[x_{m}]\) , and so is etale (V, p.~30, Prop. 3).

COROLLARY. --- For an algebraic extension E to be separable over K it is necessary and sufficient that every element of E should be a simple root of its minimal polynomial over K. It is enough to apply Prop. 5 and 6.

